\section{故障诊断实验与结果分析}
\label{sec:experiments_v5}

\subsection{引言}

前述章节分别完成了汽轮机系统健康状态建模、故障知识图谱构建以及KG-GCN故障诊断模型设计。本章在此基础上开展故障诊断算例验证。本文总体研究对象为燃煤机组汽轮机热力系统，第三章构建的系统级知识图谱覆盖主通流、再热、抽汽回热及凝汽器冷端等主要设备和参数关系；考虑不同设备的建模数据完整性和当前实验基础，本章选取超高压缸（VHP）作为代表性设备，对所提出的汽轮机系统KG-GCN故障诊断方法进行验证。

当前机组尚缺少能够覆盖多类典型故障的完整现场故障样本。为在不将构造数据表述为现场真实故障的前提下验证诊断流程，本文采用真实健康运行残差作为背景，并依据独立故障机理构造具有时间变化和空间响应差异的半物理故障场景。实验重点比较CNN、AE和KG-GCN在相同数据条件下的故障识别结果，并通过混淆矩阵分析不同模型的误诊特征。

\subsection{研究对象及故障数据集}
\label{subsec:fault_dataset_v5}

本文选取VHP作为代表性诊断对象。VHP位于主蒸汽进入汽轮机后的首个主要膨胀设备，其状态既受到主蒸汽压力、温度、流量和调节阀等入口边界影响，又通过排汽与一次再热系统相连，并通过一级抽汽支路向1号高压加热器提供抽汽。因此，在VHP局部范围内同时存在汽缸本体、抽汽支路和再热接口等不同物理区域，适合用于验证知识图谱约束下的多参数故障识别方法。

第二章VHP健康状态模型以本机实际DCS运行量作为输入，并对设备状态参数进行同步预测。完成状态点筛选后，最终选取24个状态参数用于本章故障诊断，包括VHP叶片级压力、排汽压力与温度、一级抽汽口压力与温度、一级抽汽逆止门前后温度、1号高压加热器入口抽汽压力与温度、一次低温再热接口温度以及VHP缸壁和内部蒸汽温度等。每个状态参数均保持固定DCS身份，并通过第三章设备--参数映射关系对应至VHP任务参数子图。

结合汽轮机通流部件退化、汽封泄漏机理、一级抽汽支路结构以及现有24个状态参数的可观测范围，构造正常状态和三类典型故障场景，故障类别如表~\ref{tab:fault_classes_v5}所示。

\begin{table}[htbp]
    \centering
    \caption{VHP故障诊断样本类别}
    \label{tab:fault_classes_v5}
    \begin{tabular}{p{1.0cm}p{4.1cm}p{5.3cm}p{3.4cm}}
        \toprule
        类别 & 故障场景 & 主要状态特征 & 主要作用区域 \\
        \midrule
        F0 & 正常状态 & 真实健康运行残差，不叠加故障响应 & 全部诊断节点 \\
        F1 & VHP通流部件侵蚀 & 有效通流面积及级组压力分布发生改变 & VHP本体、排汽及再热接口 \\
        F2 & 一级抽汽支路泄漏 & 抽汽支路蒸汽损失，下游抽汽状态发生衰减 & 一级抽汽及1号高加入口 \\
        F3 & VHP汽封泄漏 & 无效蒸汽泄漏引起本体及排汽热力状态改变 & VHP本体、排汽及相关接口 \\
        \bottomrule
    \end{tabular}
\end{table}

健康背景数据取自2024年6月独立运行数据。首先利用第二章GRU健康状态模型计算24个状态参数的实测值与健康预测值之差，形成真实健康残差序列。对于第$k$类故障，第$n$个半物理故障残差窗口表示为
\begin{equation}
R_{k,n}^{F}=R_n^{H}+a_{k,n}S_k\odot G_{k,n}+\epsilon_{k,n},
\end{equation}
式中：$R_n^{H}$为真实健康残差窗口；$S_k$为依据独立故障机理确定的空间响应模式；$a_{k,n}$为故障严重程度系数；$G_{k,n}$为故障随时间演化的动态函数；$\epsilon_{k,n}$为小幅随机扰动。不同样本随机改变故障发生位置和响应幅值，并采用阶跃、斜坡或Sigmoid等不同时间演化形式，以降低模型对单一固定模板的依赖。主算例中故障响应强度在训练健康残差标准差的$2.0\sim3.0\sigma$范围内随机变化，该范围仅表示清晰可观测的半物理异常强度，不对应设备实际损伤百分比。

需要说明的是，第三章用于构建主知识图谱的A0/A1关系与本章故障样本的点级空间响应模式相互独立。具体DCS测点的故障响应方向和幅值可以依据B级热力推断辅助构造，但不反向写入KG-GCN主邻接矩阵，从而避免同一组点级故障知识同时参与样本生成和分类拓扑定义造成循环验证。

为避免相邻滑动窗口跨越训练集和测试集造成信息泄漏，首先按照自然日对真实健康背景进行互斥划分，再分别在各日期内部生成半物理故障窗口。每连续5个自然日中，第1--3日作为训练集，第4日作为验证集，第5日作为测试集，且任一60~min残差窗口均不跨越自然日边界。按照上述规则，训练集、验证集和测试集分别包含1712、556和560个样本，各类别样本数量如表~\ref{tab:fault_dataset_v5}所示。

\begin{table}[htbp]
    \centering
    \caption{故障诊断数据集样本分布}
    \label{tab:fault_dataset_v5}
    \begin{tabular}{lrrrr}
        \toprule
        数据集 & F0 & F1 & F2 & F3 \\
        \midrule
        训练集 & 428 & 428 & 428 & 428 \\
        验证集 & 139 & 139 & 139 & 139 \\
        测试集 & 140 & 140 & 140 & 140 \\
        \bottomrule
    \end{tabular}
\end{table}

\subsection{实验设置与评估指标}
\label{subsec:experiment_setting_v5}

为验证本文所提KG-GCN模型的故障诊断性能，选取CNN和AE作为对比算法。CNN将$24\times60$的多参数残差窗口作为二维输入，通过局部卷积核提取参数--时间特征；AE首先利用无监督重构学习低维潜在表示，再通过分类层完成故障识别；KG-GCN则将24个状态残差序列加载至知识图谱对应参数节点，并依据第三章得到的参数关联邻接矩阵进行图卷积特征传播。三种模型采用完全相同的训练、验证和测试样本。

KG-GCN参数节点数量为24，残差窗口长度为60。图卷积编码器包含2层GCN，每个节点首先映射为32维隐藏特征，并沿知识图谱邻接关系完成两层信息聚合。无标签预训练阶段随机遮蔽15\%的残差位置，通过恢复被遮蔽数据训练图卷积编码器；预训练完成后固定GCN编码器参数，仅利用故障标签训练全连接分类层。分类层隐藏维数为96，优化器采用AdamW，权重衰减为$1\times10^{-4}$。主要实验参数如表~\ref{tab:kg_setting_v5}所示。

\begin{table}[htbp]
    \centering
    \caption{KG-GCN模型主要实验参数}
    \label{tab:kg_setting_v5}
    \begin{tabular}{ll}
        \toprule
        参数 & 取值 \\
        \midrule
        参数节点数量 & 24 \\
        残差窗口长度 & 60 \\
        GCN隐藏层数量 & 2 \\
        节点隐藏特征维数 & 32 \\
        Mask比例 & 15\% \\
        分类层隐藏维数 & 96 \\
        优化器 & AdamW \\
        权重衰减 & $1\times10^{-4}$ \\
        \bottomrule
    \end{tabular}
\end{table}

本文采用准确率（Accuracy）、精确率（Precision）、召回率（Recall）和F1值评价不同算法的故障诊断性能，并利用混淆矩阵进一步分析各故障类别的识别结果。对于第$k$类故障，精确率、召回率和F1分别定义为
\begin{equation}
P_k=\frac{TP_k}{TP_k+FP_k},
\end{equation}
\begin{equation}
R_k=\frac{TP_k}{TP_k+FN_k},
\end{equation}
\begin{equation}
F1_k=\frac{2P_kR_k}{P_k+R_k},
\end{equation}
式中：$TP_k$、$FP_k$和$FN_k$分别表示第$k$类故障的正确分类、误报和漏报样本数量。由于本章属于多类别诊断任务，对各类别指标进行宏平均，以避免某一类别样本对综合结果产生过强影响：
\begin{equation}
Precision_{macro}=\frac{1}{C}\sum_{k=1}^{C}P_k,
\end{equation}
\begin{equation}
Recall_{macro}=\frac{1}{C}\sum_{k=1}^{C}R_k,
\end{equation}
\begin{equation}
F1_{macro}=\frac{1}{C}\sum_{k=1}^{C}F1_k .
\end{equation}
整体准确率定义为
\begin{equation}
Accuracy=\frac{N_{correct}}{N_{all}},
\end{equation}
其中$N_{correct}$和$N_{all}$分别为正确分类样本数和测试样本总数。

\subsection{诊断结果与对比}
\label{subsec:results_v5}

表~\ref{tab:main_result_v5}给出了CNN、AE和KG-GCN三种算法在同一测试集上的诊断结果。为保持工程算例结果的直观性，表中首先给出固定随机种子123条件下的主实验结果。可以看出，KG-GCN的诊断准确率达到98.57\%，高于CNN的97.86\%和AE的95.54\%；其宏平均精确率、召回率和F1分别达到98.60\%、98.57\%和98.58\%，四项指标均为三种模型中最高。CNN同样能够获得较高整体识别率，但在部分具有相似局部残差特征的故障类别之间仍存在误判；AE整体结果低于另外两种方法。

\begin{table}[htbp]
    \centering
    \caption{CNN、AE和KG-GCN算法故障诊断结果对比}
    \label{tab:main_result_v5}
    \begin{tabular}{lcccc}
        \toprule
        模型算法 & Accuracy/\% & Precision/\% & Recall/\% & F1/\% \\
        \midrule
        CNN & 97.86 & 97.89 & 97.86 & 97.85 \\
        AE & 95.54 & 95.57 & 95.54 & 95.54 \\
        \textbf{KG-GCN} & \textbf{98.57} & \textbf{98.60} & \textbf{98.57} & \textbf{98.58} \\
        \bottomrule
    \end{tabular}
\end{table}

为进一步分析不同模型的具体误诊类别，在相同测试数据和随机种子条件下分别统计CNN、AE和KG-GCN的混淆矩阵。矩阵的行表示真实故障类别，列表示模型预测类别，三种模型结果分别为
\begin{equation}
C_{CNN}=\begin{bmatrix}
139&1&0&0\\
5&135&0&0\\
0&2&134&4\\
0&0&0&140
\end{bmatrix},
\end{equation}
\begin{equation}
C_{AE}=\begin{bmatrix}
134&1&1&4\\
6&130&4&0\\
3&3&134&0\\
2&1&0&137
\end{bmatrix},
\end{equation}
\begin{equation}
C_{KG-GCN}=\begin{bmatrix}
139&1&0&0\\
4&136&0&0\\
1&1&138&0\\
1&0&0&139
\end{bmatrix}.
\end{equation}

从CNN混淆矩阵可以看出，测试集中共有12个样本被错误分类。其中F2一级抽汽支路泄漏有4个样本被识别为F3 VHP汽封泄漏，是CNN中较明显的一类交叉误判。CNN主要利用局部卷积核提取参数--时间矩阵中的邻域特征，当两类故障同时在抽汽或VHP排汽附近形成幅值和短时变化相近的残差时，仅依赖局部数据模式难以显式区分异常究竟沿一级抽汽支路传播，还是主要发生于VHP本体及汽封相关区域，因此容易在共享局部特征的故障之间产生混淆。

AE测试集中共有25个样本被错误分类，误判数量高于CNN和KG-GCN，且错误分布更加分散。F0正常状态、F1通流部件侵蚀和F2一级抽汽支路泄漏均出现向多个类别误判的情况。AE的训练目标首先侧重于压缩和重构输入残差，当不同故障样本在整体残差幅值或潜在空间表示上具有相似性时，其低维特征可能落入较接近的区域。由于AE没有显式引入参数之间的设备归属和热力连接关系，潜变量表示难以进一步利用异常所处物理位置区分具有相似全局响应的故障类别。

相比之下，KG-GCN混淆矩阵中的样本分布最集中于主对角线，560个测试样本中仅有8个发生误诊。F2一级抽汽支路泄漏仅有2个样本被错误分类，F3 VHP汽封泄漏仅有1个样本被识别为正常状态。KG-GCN将第三章知识图谱中的参数连接关系转化为特征传播路径：一级抽汽口、逆止门前后及1号高压加热器入口构成连续抽汽支路，VHP本体压力、排汽状态和一次再热接口则形成另一组具有明确设备位置的关联节点。图卷积沿这些实际结构关系逐层聚合残差特征，使模型不仅利用单个参数的异常幅值和时间变化，还能够利用异常在设备网络中的空间位置及邻域协同关系。因此，当不同故障共享部分监测参数时，KG-GCN仍能够借助参数之间的结构关系降低局部特征相似造成的误判。

为检验上述结果是否依赖单次随机初始化，在保持数据集和模型超参数不变的情况下采用3个随机种子重复训练。CNN、AE和KG-GCN的平均Accuracy分别为98.33\%、96.43\%和98.45\%，平均Macro-F1分别为98.33\%、96.43\%和98.46\%。其中KG-GCN的Macro-F1标准差仅为0.11\%，低于CNN的0.45\%和AE的0.82\%，说明所提模型在不同随机初始化条件下具有较稳定的诊断结果。该重复实验仅作为主算例的稳定性核验，不改变前述固定随机种子下的主要对比结论。

考虑实际电站中带有明确故障标签的数据通常较少，进一步对第四章掩码预训练策略进行补充验证。当仅使用20\%带标签窗口训练故障分类阶段时，直接随机初始化训练KG-GCN的Macro-F1为94.48\%$\pm$0.93\%；采用无标签残差掩码预训练并在后续固定图卷积编码器后，Macro-F1提高至96.67\%$\pm$0.27\%，平均提升2.19个百分点。由此可见，在故障标签数量降低时，利用未标注残差提前学习知识图谱约束下的参数关联能够为分类阶段提供更加稳定的结构化特征表示。

\subsection{本章小结}

本章选取VHP作为代表性设备，对本文提出的汽轮机系统KG-GCN故障诊断方法进行了算例验证。首先基于2024年6月真实健康运行残差和独立故障机理构造正常状态、VHP通流部件侵蚀、一级抽汽支路泄漏以及VHP汽封泄漏四类诊断样本，并按照自然日完成训练、验证和测试数据划分；随后以CNN和AE作为对比算法，采用Accuracy、Precision、Recall、F1和混淆矩阵评价诊断结果。固定随机种子主实验中，KG-GCN的Accuracy和F1分别达到98.57\%和98.58\%，均高于CNN和AE，且混淆矩阵中错误样本最少。进一步的重复实验表明KG-GCN具有较低的随机初始化波动；当仅使用20\%带标签故障窗口时，掩码预训练使Macro-F1平均提高2.19个百分点。上述结果说明，将健康状态残差映射至汽轮机知识图谱参数节点，并利用图卷积沿实际参数关系传播特征，可以提高具有多参数耦合特征的故障分类效果。

需要指出的是，本章诊断样本属于基于真实健康残差与独立故障机理构造的半物理故障场景，并非本文机组现场真实故障记录。因此，本章结论用于验证所提方法在机理一致构造场景下的有效性，后续仍需结合现场故障案例或经校准的热力仿真结果进一步检验工程泛化能力。